% HD-S35. GPT-6 Astra Pro, /root/sixfolds; inherited effort not exposed.
% 21 November 2026. All rank-four Gaussian finite coefficient groups.
\subsection{HD-S35: all rank-four Gaussian finite sources}
\label{sec:rank-four-finite-scope}

Let \(G\subset GL_4(K)\), \(K=\mathbb Q(i)\), be a finite
absolutely irreducible actual matrix group with character field exactly
\(K\). Use its natural primitive rational group-algebra factor of the
Jacobian of a connected full-\(G\) sphere cover, with the Gaussian corner
and restricted theta polarization specified as in HD-S23. Every such
Weil sixfold of signature \((3,3)\) in a period-rank-nine family has
ordinary Hermitian determinant \([-1]\) or \([-3]\).
The group need not be generated by reflections or be a scalar extension
of a reflection group. The proof supplies the normal-subgroup reduction,
all remaining normalizers, and the actual primitive-factor map.

\paragraph{1. The forced reflection subgroup and its actual blocks.}
Write \(g_1,\ldots,g_b\) for the nonidentity local monodromies and
\(c_j=\operatorname{codim}_K V^{g_j}\). The parabolic Euler formula,
with both global invariant spaces zero by irreducibility, is
\[
               6=-2\cdot4+\sum_jc_j,
               \qquad \sum_jc_j=14.                         \tag{S35.1}
\]
For fixed finite monodromy data, the Hurwitz map to the configuration
of the \(b\) branch values is finite: the fibre is a subset of the
finite set of generating product-one tuples in \(G\), modulo
simultaneous conjugacy. Quotienting configurations by \(PGL_2\)
gives dimension \(b-3\). Thus period rank nine forces \(b\ge12\).
Every \(c_j\ge1\), so
\(\sum_j(c_j-1)=14-b\le2\). At least ten entries are actual
reflections.

Let \(W\) be the subgroup of \(G\) generated by all its reflections.
Conjugation preserves rank \((s-1)\), so \(W\) is a nontrivial
normal subgroup. Average a positive \(K\)-Hermitian form over \(G\).
Every root line \(\operatorname{im}(s-1)\) is \(K\)-defined,
and its orthogonal complement is \(\ker(s-1)\). Join two such
lines when their Hermitian pairing is nonzero. Spans of distinct
connected components are orthogonal \(K\)-subspaces. Their sum is
nonzero and \(G\)-invariant, hence is all \(V\).

Each component span is an irreducible reflection module. Indeed an
invariant complex subspace and its orthogonal complement would place
every rank-one image wholly in one of them; the nonorthogonality
connections then place the whole connected component in one. The
normality of \(W\) makes \(G\) permute these spans, and irreducibility
makes that action transitive. Their common rank consequently divides
four. Rank-one components give the literal \(K\)-defined line system
of HD-S34, whose determinant is \([-1]\). Rank-two components give
exactly the induced block system and quotient-cover Shapiro map of
HD-S30; its period rank is at most seven. The remaining case is that
\(W\) itself is irreducible of rank four.

The exact classification and trace restrictions in HD-G25 leave
\[
 A_4,\ B_4=G(2,1,4),\ D_4=G(2,2,4),\ F_4,
 \quad G(4,p,4)\ (p=1,2,4),\quad G_{29},G_{31}.             \tag{S35.2}
\]
Here the restriction of the given representation to \(W\) has its
traces in \(K\), which is exactly what those trace tests use.
The finite intertwiner system in HD-G25 gives an actual \(K\)-matrix
carrying this restriction to the indicated standard root model.
Its pullback form is a positive rational multiple of the original
form. On parabolic cochains this is the literal coordinatewise map;
on primitive Jacobians its lattice-cleared isogeny has the exact
scalar \(M^2a c_0/c_{\rm given}\) of (G25.3a). Its sixth power
is a norm. Thus we work with these specified models and maps.

\paragraph{2. The monomial normalizers recover the coordinate lines.}
All the following normalizers preserve the standard Hermitian form
up to its common positive scalar, by absolute irreducibility of \(W\).
For \(B_4\), its coordinate reflections form a conjugacy class of
size four, whereas its other reflections form one of size twelve.
A normalizer permutes conjugacy classes and preserves their sizes,
so it preserves the four coordinate root lines.

In \(G(4,1,4)\), the roots of the order-four reflections are precisely
those same four lines. In \(G(4,2,4)\), the coordinate reflections
form a class of size four and the other reflections a class of size
twenty-four. To check the latter class assertion, permutations move
the unordered pair of coordinates arbitrarily, and diagonal conjugation
changes its phase by any element of \(\mu_4\); a third coordinate
compensates the exponent so the diagonal matrix remains in the group.
The coordinate class is separately transitive by permutations.
These explicit conjugations prove the stated sizes and distinguish
the coordinate lines intrinsically.

For \(G(4,4,4)\), its twenty-four root lines are
\[
           K(e_j-\zeta e_k),\qquad j<k,\quad\zeta\in\mu_4.
\]
Give the lines the graph relation that their normalized inner-product
absolute value is \(1/\sqrt2\). Two distinct lines have this relation
exactly when they have the same pair \(\{j,k\}\) and their phases
have ratio \(i\) or \(-i\). Different coordinate pairs sharing one
coordinate have normalized absolute value \(1/2\), and disjoint
pairs have value zero. The graph therefore has six connected
four-cycles. Their spans are exactly the six planes
\(Ke_j+Ke_k\). The nonzero intersections of three distinct such
planes are exactly the four coordinate lines. Every normalizer
preserves the graph, its component spans, and these intersections.
Consequently it also preserves the coordinate line system.

In every case in this paragraph, the whole \(G\) is thus monomial
on four actual \(K\)-lines. Apply HD-S34, including its induced
stabilizer character, integral local and global maps, noncentral
corner action, and cyclic quotient realization. The determinant
is \([-1]\).

\paragraph{3. The \(A_4\) and nonreal primitive normalizers.}
For \(A_4\), use the augmentation space of \(K^5\), with root
lines \(K(e_j-e_k)\). Two distinct such lines have nonzero inner
product precisely when the corresponding edges of the complete graph
on five vertices share a vertex. The four-cliques of this line graph
are exactly the five stars: pairwise intersecting four two-element
subsets must have a common element, since a triangle admits no fourth
edge intersecting all three. A normalizer permutes these five stars
and hence induces an actual permutation in \(S_5=W\) on every
root line. After that permutation is removed, the unitary map fixes
every root line. Its scalar eigenvalues agree on nonorthogonal lines,
so connectivity makes it scalar on the whole augmentation space.
Thus every normalizer element is \(aw\), \(w\in W\).
Its entries and those of \(w\) lie in \(K\), so \(a\in K\);
finite order gives \(a\in\mu_4\). Since \(W\subset G\), each
such scalar separately belongs to \(G\). Character field \(K\)
therefore forces \(iI\in G\), and
\(G=\mu_4\times S_5\) in this representation. This is the actual
\(N=5\) source of HD-S33, with determinant \([-1]\) at rank nine.

For \(G_{29}\) and \(G_{31}\), use Cohen's corrected normalizer
proposition, \emph{Erratum: Finite complex reflection groups},
Ann. Sci. ENS 11 (1978), printed p.613. It replaces (5.14) of the
1976 paper: for a nonreal primitive reflection group \(W\) normal
in a finite unitary group \(G\), one has \(G\subset\mu_\infty W\),
apart from the stated rank-three \(W(L_3)\) exception. The original
scan explicitly has the normal-subgroup symbol; this hypothesis is
provided here by the construction of \(W\). The primary scan is
\url{https://www.numdam.org/item/ASENS_1978_4_11_4_613_0.pdf},
PDF page two. Both rank-four groups have literal Gaussian matrices
and contain \(iI\), as rebuilt in HD-S32. Writing \(g=aw\) makes
\(a\in K\cap\mu_\infty=\mu_4\subset W\). Hence \(G=W\)
in these cases, and HD-S32 gives \([-1]\) or \([-3]\).

\paragraph{4. Exact real and Gaussian normalizers of \(D_4,F_4\).}
Use \(I_4\) as coefficient Gram matrix. Put
\[
\begin{split}
 \mathcal S&=\{\pm e_j\}\ \cup\
       \{\tfrac12(\pm e_1\pm e_2\pm e_3\pm e_4)\},\\
 \mathcal L&=\{\pm e_j\pm e_k:j<k\}.
\end{split}
\]
Both sets have twenty-four elements, of squared norms one and two.
The long roots \(\mathcal L\) give \(D_4\); their union with
\(\mathcal S\) gives \(F_4\). The four generators and exact
order \(|W(F_4)|=1152\) were rebuilt in HD-G25. They preserve
\(\mathcal S\) and \(\mathcal L\) separately. The exact root
certificate accompanying this proof also reconstructs
\(|W(D_4)|=192\) and its normality in \(W(F_4)\).

First pass from a unitary complex normalizer to a real one through
an explicit scalar. If \(g\) normalizes a real absolutely irreducible
\(W\), then \(g^{-1}\bar g\) commutes with every element of
\(W\), hence is scalar. All nonzero entries of \(g\) therefore
have the same complex phase up to a real sign. Choose that phase
\(a\) with \(|a|=1\); then \(g=aR\) with \(R\) real orthogonal
and normalizing \(W\). The real normalizer is finite because it
permutes a finite spanning root-line set and its kernel consists of
\(\pm I\). Thus \(R\) has finite order, and finite order of
\(g\) makes \(a\) a root of unity.

An orthogonal normalizer of \(D_4\) permutes the actual long-root
set. An ordered orthogonal four-frame in that set has respectively
\(24,6,4,2\) choices for its successive vectors. These counts follow
by taking the first vector to \(e_1+e_2\), then listing its six
orthogonal roots; the subsequent two orthogonality equations leave
four and two choices. Equivalently the attached exact enumeration
checks every frame, not only one orbit. An orthogonal map is determined
by a frame. There are at most \(24\cdot6\cdot4\cdot2=1152\)
such maps. The subgroup \(W(F_4)\) already attains that order and
preserves \(\mathcal L\). Therefore
\[
               N_{O_4(\mathbb R)}(W(D_4))=W(F_4).            \tag{S35.3}
\]
All its matrices are rational. Thus a finite Gaussian normalizer
has its phase \(a\in K\cap\mu_\infty=\mu_4\), giving the
actual containing matrix group \(\mu_4W(F_4)\), of order 2304.

For \(F_4\), consider its twenty-four normalized root lines. Join
two when their inner-product absolute value is \(1/\sqrt2\).
The displayed coordinates show a connected bipartite graph with
parts the twelve short and twelve long lines: each vertex has six
neighbours in the opposite part, and successive shared-coordinate
edges connect all vertices. The exact finite certificate verifies
this same graph directly. Any normalizer either preserves or exchanges
its two parts. A part-preserving orthogonal map permutes
\(\mathcal S\); its ordered orthogonal four-frames again have
successive counts \(24,6,4,2\). Thus the part-preserving group
has order at most 1152 and is exactly \(W(F_4)\).

Set
\[
 M=\begin{pmatrix}1&1&0&0\\1&-1&0&0\\0&0&1&1\\0&0&1&-1\end{pmatrix},
 \qquad u=M/\sqrt2.
\]
It is an orthogonal involution. The exact identities
\(M\mathcal S=\mathcal L\) and
\(\tfrac12M\mathcal L=\mathcal S\) show that \(u\) exchanges
the normalized root parts. Multiplying a part-exchanging map by
\(u\) makes it part-preserving. Consequently the real normalizer
is exactly \(W(F_4)\cup uW(F_4)\).

For a Gaussian \(g=aR\), the coset \(R\in W(F_4)\) again gives
\(a\in\mu_4\). In the other coset \(R=uw\), the nonzero rational
entries of \(Mw\) imply \(a/\sqrt2\in K\). Therefore \(a^2\)
is a root of unity in \(K\), hence belongs to \(\mu_4\).
The possibilities \(a^2=1,-1\) would put \(a\in K\), contradicting
\(a/\sqrt2\in K\); thus \(a\) is a primitive eighth root.
It follows literally that
\[
 G\ \subset\ \mathcal N=\langle W(F_4),U\rangle,
 \qquad U=\frac{1+i}{2}M,\qquad U^2=iI_4.                  \tag{S35.4}
\]
The four cosets \(W,UW,iW,iUW\) are distinct and exhaust
\(\mathcal N\), so \(|\mathcal N|=4608=2^9 3^2\).
All their entries lie in \(\tfrac12\mathbb Z[i]\).
Indeed columns of a matrix in \(W(F_4)\) belong to \(\mathcal S\);
applying \(M\) gives integral columns, so the \(UW\) coset has
the asserted half-Gaussian entries. These actual matrix groups,
not their abstract orders alone, provide the next integral lattice.

\paragraph{5. Perfect parabolic lattice and its determinant.}
For either remaining case put
\[
          R=\mathbb Z[i,1/6],\quad L=R^4,
          \quad h_0(x,y)=\bar x^t y.                        \tag{S35.5}
\]
The actual groups just obtained preserve this explicit self-dual
lattice. Every local monodromy has order with prime factors only
2 and3. For its order \(o\), the two integral endomorphisms
\[
 e_g=o^{-1}\sum_{j=0}^{o-1}g^j,
 \qquad B_g=o^{-1}\sum_{j=0}^{o-1}(o-1-j)g^j
\]
satisfy \((1-g)B_g=1-e_g\). Thus both the invariant module and
the coinvariant module are the free summand \(e_gL\).
Both groups contain \(-I\in W(D_4)\); since2 is invertible,
absolute zeroth homology for \(L\) and its dual vanishes.

Apply the actual integral Poincar\'e--Lefschetz and primitive-image
maps proved in HD-S27 and spelled out again in HD-S32. The image
of absolute first homology in relative first homology is primitive
because its cokernel is the free boundary zeroth homology. The two
absolute-to-relative maps are adjoint under perfect cross-pairings;
extending a functional from the primitive image gives the perfect
parabolic pairing. This identifies with the middle extension through
the local summands \(e_gL\). Hence the rank-six skew-Hermitian
cup matrix \(S\) is perfect over \(R\), as is \(H=iS\).
Its determinant belongs to \(R^\times\cap\mathbb Q\), and
signature \((3,3)\) makes it negative. Thus
\[
                   \det H=-2^a3^b,
                   \qquad a,b\in\mathbb Z.
\]
Since2 is the Gaussian norm of \(1+i\), and every even power
of3 is a norm, the ordinary class is \([-1]\) or \([-3]\).

\paragraph{6. The actual primitive corner, including a noncentral Gaussian action.}
This paragraph retains the abelian-variety map even when the original
\(G\) contains no scalar \(iI\). It applies to the two groups of
paragraphs4--5. Put \(v=(1+i)e_1\), so \(h_0(v,v)=2\), and
\[
 e=\frac4{|G|}\sum_{g\in G}
       \operatorname{Re}h_0(v,\rho(g^{-1})v)\,g,
       \qquad D=|G|/4.                                     \tag{S35.6}
\]
The order of \(W\) divides \(|G|\), so \(D\) is an integer.
The half-Gaussian entry bound makes \(De\in\mathbb Z[G]\).
Schur orthogonality gives the rank-one orthogonal projector onto
\(Kv\) on \(\rho\), its conjugate projector on \(\bar\rho\),
and zero on every other complex irreducible. Thus \(e\) is a
primitive rational idempotent with corner \(K\).
With left deck maps \(L_g\) and actual Jacobian action
\(a_g=L_{g,*}\), let \(B=\operatorname{im}(a_{De})\).

Choose a conjugacy class \(C\) whose character value is nonreal;
one exists because the character field is exactly \(K\). Its central
class sum acts on \(\rho\) by
\[
          \lambda=\frac{|C|}{4}\operatorname{tr}\rho(c)
                         =a+bi,\quad a,b\in\mathbb Q,
                         \quad b\ne0.
\]
Let \(A_C=\sum_{c\in C}a_c\) on \(B\). The Gaussian action on
this corner is the literal rational endomorphism
\[
                        \iota_B=(a-A_C)/b.                  \tag{S35.7}
\]
To verify the sign, the cohomology action is
\(T_g=(L_g^{-1})^*\), and the matrix coefficient
\(j_{v,x}(g)=h_0(v,\rho(g)x)/2\) is conjugate-linear in \(v\).
Hence \(T_Cj_{v,x}=\bar\lambda j_{v,x}\), and Poincar\'e duality
identifies \(A_C\Phi(x)=\Phi(\bar\lambda x)\) for
\(\Phi=\operatorname{PD}(j+\bar j)\). Therefore
\(\iota_B\Phi(x)=\Phi(ix)\), proving \(K\)-linearity,
\(\iota_B^2=-1\), and the inverse-deck convention. The Rosati
adjoint of a deck action is its inverse; the inverse class has scalar
\(\bar\lambda\). It follows that \(\iota_B^\dagger=-\iota_B\).

The same real Schur computation as in HD-S23 now gives the exact
polarized identity
\[
                   h_B(\Phi x,\Phi y)=D\,H(x,y).            \tag{S35.8}
\]
It is the sum of the two conjugate cross terms:
\(2|G|/(4h_0(v,v))=D\). The rational Hodge map \(\Phi\),
with an integer clearing its homology denominators, is an actual
isogeny of the indicated complex tori and abelian factors. Thus the
ordinary determinant of the restricted theta form on \(B\) differs
from that of \(H\) by the norm \(D^6=N(D^3)\) and a basis norm.
This establishes the assertion for the actual source polarization.

If an integral \(\mathbb Z[i]\)-action is desired, there is an
explicit lattice map. For \(B=V_\mathbb R/\Lambda\), choose
\(d>0\) with \(d\iota_B\Lambda\subset\Lambda\), put
\(\Lambda'=\Lambda+\iota_B\Lambda\), and give \(V_\mathbb R/\Lambda'\)
the positive Riemann form \(d^2E\). It is integral because
\(d\Lambda'\subset\Lambda\), and \(\Lambda'\) is \(\iota_B\)-stable.
The quotient isogeny pulls this form back to \(d^2E\); its
rank-six determinant change is again a norm. This specifies the
integral order transport without assuming the class-sum action itself
was integral.

Different primitive coefficient lines are related by the actual
nonzero matrix-coefficient corner map, cleared of denominators, as
in HD-S34. For rank-nine families, compatible flat changes of
polarization under a fixed isogeny have the rational scalar comparison
proved in HD-S17. An originally nonfaithful finite representation is
also reduced by an actual map: quotient the covering curve by its
kernel, then use pullback and norm, whose composite is multiplication
by the kernel order. This multiplies the Jacobian polarization by
that positive integer and the rank-six determinant by a norm.
These maps justify exactly the source and polarization scope claimed.

\paragraph{Whole-proof check and finite supplement.}
The author read the complete proof, including every normalizer and
the class-sum sign. The exact script
\path{rank_four_normalizer_certificate.py} rebuilds the real
root groups, their Gaussian normalizer, all ordered orthogonal frames,
the connected bipartite root graph, the six four-cycles of
\(G(4,4,4)\), and the five \(A_4\) stars. It verifies orders 192,
1152 and 4608, half-Gaussian matrices, \(U^2=iI\), and normality.
Those computations support the explicit finite geometry; the
universal period, block, parabolic, and Jacobian arguments are the
proof above. No dominance or supplier existence is inferred from a
dimension count or from the allowed determinant list.
